\\ 
\textit{Figures.} Please see the generated plots in Figure~\ref{fig:gain_study}. 

\textit{Gain discussion.} $k_v$ has the larger effect of the two gains on
performance, determining whether the system converges without overshooting and
having to turn back, where the model difference is much more significant. If the
gain is too low, the controller does not keep up with the desired velocity and
performs poorly near the final position, not slowing down fast enough. If the
gain is high enough, however, there is not substantial difference in
performance, only varying slightly in transient characteristics. 

$k_\theta$ determines how direct the path is the final goal, with higher gains
causing more initial overshoot (likely partly due to model mismatch), but having
better final error. The lower gains perform better transiently but exhibit
more steady-state error at the final orientation.

There is no gain pair that performs the best across all of the errors (for both
transient and final), so the choice should be made on desired characteristics.
Likely, the most important error is position, so the gain pair $(k_v, k_\theta)
= (2.5,3.5)$ would be a good choice, causing more initial overshoot in
orientation but having good overall final error characteristics.

\begin{figure}[H]
    \centering
    \begin{subfigure}[b]{0.48\linewidth}
        \centering
        \includegraphics[height=5.35cm]{../../results/gain_study/position_error.png}
    \end{subfigure}
    \hfill
    \begin{subfigure}[b]{0.48\linewidth}
        \centering
        \includegraphics[height=5.35cm]{../../results/gain_study/orientation_error.png}
    \end{subfigure}
    \captionsetup{font={color=blue}}
    \caption{Gain study plots.}
    \label{fig:gain_study}
\end{figure}